\documentclass[11pt]{article}
\usepackage{fontspec}
\usepackage[french,english]{babel}
\usepackage{amsmath}
\usepackage{amssymb}
\usepackage{amsthm}
\usepackage{mathrsfs}
\usepackage{tikz-cd}
\usepackage{enumitem}
\usetikzlibrary{babel}

\newcommand{\printedpage}[1]{\marginpar{\footnotesize p.~#1}}
\newcommand{\msp}[1]{\marginpar{\footnotesize [#1]}}
\newcommand{\uFl}{\underline{\mathrm{Fl}}}
\newcommand{\Fl}{\mathrm{Fl}}
\newcommand{\Filt}{\mathrm{Filt}}
\newcommand{\cM}{\mathcal{M}}
\newcommand{\card}{\operatorname{card}}
\newcommand{\circled}[1]{\tikz[baseline=(c.base)]\node[draw,circle,inner sep=0.5pt](c){#1};}
% footnote mark printed as superscript "(n)"
\newcommand{\pfn}[2]{\textsuperscript{(}\footnote[#1]{#2}\textsuperscript{)}}

\begin{document}

\printedpage{9}
% head: A. Grothendieck : Dérivateurs, Ch. XV, Édition des Universités de Montpellier II et Paris VII
% foot: Version 5 novembre 2010

{\itshape whence by iteration an ordinal direct system of endofunctors $f\longmapsto p_\alpha(f)$ of $\uFl(\cM)$, and of homomorphisms $i_\alpha : \mathrm{id}_{\uFl(\cM)} \longrightarrow p_\alpha$, \textup{i.e.}}
\[
\begin{tikzcd}[column sep=small]
X \arrow[rr, "i_\alpha(f)"] \arrow[dr, "f"'] & & X_\alpha(f) \arrow[dl, "p_\alpha(f)"] \\
& Y\,, &
\end{tikzcd}
\]
{\itshape forming an ordinal direct system of factorizations of $f$ functorial in $f$.}

\msp{page 11}

{\itshape This being so, we have:}
\begin{enumerate}[label={}, leftmargin=*]
\item[a)] {\itshape For every ordinal $\alpha$, we have $i_\alpha(f)\in\Phi$.}
\item[b)] {\itshape Suppose that the sources of the $i\in C_0$ are accessible and then let $\pi$ be a cardinal in $\mathfrak{U}$ such that they are $\pi$-accessible. Then for every ordinal number $\alpha$ large relative to $\pi$, we have $p_\alpha(f)\in C_{0*}$, hence}
\[
f \;=\; p_\alpha(f)\, i_\alpha(f) \qquad\quad \textit{with} \qquad\quad i_\alpha(f)\in\Phi\,, \qquad\quad p_\alpha(f)\in\Psi\,.
\]
\end{enumerate}

We have a slightly more precise result, as follows:

\textbf{Corollary 1.} \textit{Let $\Phi_0$ be a subset of} $\Fl(\cM)$ \textit{such that}
\[
C_0''' \;\subset\; \Phi_0 \;\subset\; \Phi \qquad\quad (\textit{hence}\ \ C_0\subset\Phi_0\subset\Phi)\,,
\tag*{(1.24)}
\]
\textit{where $C_0'''\subset\Phi$ is defined in} 1.19, \textit{p.} 8. \textit{For example, if $\Phi_0$ is a subset of} $\Fl(\cM)$ \textit{satisfying}
\[
C_0 \;\subset\; \Phi_0 \;\subset\; \Phi\,,
\]
\textit{and stable under cobase change, under small amalgamated sums} (\textit{cf.} (1.18), p.~8), \textit{under strict ordinal $\varinjlim$} (\textit{cf.} (1.19), p.~8), \textit{it satisfies} (1.24). (\textbf{N.B.} \textit{If $\cM$ admits simple amalgamated sums, such $\Phi_0$ exist trivially, it suffices to take $\Phi_0=\Phi$.})

\msp{page 12}

\textit{This being so, for every arrow $f$ in $\cM$, and every ordinal $\alpha$, we have}
\[
i_\alpha(f) \;\in\; \Phi_0\,.
\]

\textbf{Corollary 2.} \textit{Suppose that the sources of the arrows $i\in C_0$ are accessible, and let $\Phi_0$ be as in Corollary} 1 (e.g. $\Phi_0=C_0'''$). \textit{Then the pair $(\Phi_0,\Psi)$ is a Quillen precouple,} i.e.
\begin{enumerate}[label={}, leftmargin=*]
\item[a)] $\Phi_0\longleftrightarrow\Psi$, i.e. $\Psi\subset\Phi_{0*}$ (\textit{or what amounts to the same thing, $\Phi_0\subset\Psi^*$}), \textit{and}
\item[b)] \textit{we have the Quillen factorization property.}
\end{enumerate}
\textit{This therefore implies that $\Phi$ $(=\widetilde{C_0}=\widetilde{\Phi_0})$ consists of the arrows which are direct factors of arrows of $\Phi_0$. Thus if $\Phi_0$ is stable under direct factors, we have $\Phi_0=\Phi$,} i.e. \textit{$\Phi_0$ is left closed.}

\printedpage{10}
% head: A. Grothendieck : Dérivateurs, Ch. XV, Édition des Universités de Montpellier II et Paris VII
% foot: Version 5 novembre 2010

\section*{2\quad Variation of the factorizations $f=p_\alpha i_\alpha$ with the generating subset $C_0$}

I say that the ordinal direct system of functors $p_\alpha=p_\alpha^{C_0} : \uFl(\cM)\longrightarrow\uFl(\cM)$,
\msp{page 13}
equipped with the arrows
\[
i_\alpha^{C_0} \;:\; \mathrm{id}_{\uFl(\cM)} \longrightarrow p_\alpha^{C_0}\,,
\]
depends functorially on $C_0$, when one regards $C_0$ as an element of the ordered set of \textit{small} subsets of $\Fl(\cM)$, consisting of cosquarable arrows (thus of the ordered set
\[
\mathfrak{P}_{\mathfrak{U}\text{-small}}(\Fl_{\text{cosquar.}}(\cM)) \quad ).
\]
If we have an inclusion
\[
C_0 \;\hookrightarrow\; C_1\,,
\tag*{(2.1)}
\]
we deduce from it for every ordinal $\alpha$ a functorial homomorphism
\[
u_\alpha \;=\; u_\alpha^{C_1,C_0} \;:\; X_\alpha^{C_0} \;\longrightarrow\; X_\alpha^{C_1}\,,
\tag*{(2.2)}
\]
making commutative the squares deduced from $f : X\longrightarrow Y$ in $\cM$,
\[
\begin{tikzcd}
X \arrow[r, "i_\alpha^{C_0}(f)"] \arrow[d, "\mathrm{id}"'] & X_\alpha^{C_0}(f) \arrow[r, "p_\alpha^{C_0}(f)"] \arrow[d, "u_\alpha(f)"] & Y \arrow[d, "\mathrm{id}"] \\
X \arrow[r, "i_\alpha^{C_1}(f)"] & X_\alpha^{C_1}(f) \arrow[r, "p_\alpha^{C_1}(f)"] & Y\,,
\end{tikzcd}
\tag*{(2.3)}
\]
and compatible with the transition morphisms for variable $\alpha$,
\[
\begin{tikzcd}
X_\alpha^{C_0} \arrow[r, "{\mathrm{trans}_{\beta,\alpha}^{C_0}}"] \arrow[d, "u_\alpha"'] & X_\beta^{C_0} \arrow[d, "u_\beta"] \\
X_\alpha^{C_1} \arrow[r, "{\mathrm{trans}_{\beta,\alpha}^{C_1}}"] & X_\beta^{C_1}\,.
\end{tikzcd}
\tag*{(2.4)}
\]

\msp{page 14}

To define the $u_\alpha$, satisfying the preceding conditions, it suffices to define them for $\alpha=1$ ``and to iterate'':
\[
\begin{tikzcd}[column sep=small]
& X_1^{C_0}(f) \arrow[dd, "u_1(f)"] \arrow[dr, "p_1^{C_0}(f)"] & \\
X \arrow[ur, "i_1^{C_0}(f)"] \arrow[dr, "i_1^{C_1}(f)"'] & & Y \\
& X_1^{C_1}(f) \arrow[ur, "p_1^{C_1}(f)"'] & {}.
\end{tikzcd}
\tag*{(2.5)}
\]
Now $X_1^{C_0}(f)$ is defined as a direct sum in $X\backslash\cM/Y$, relative to an index set $\Delta(C_0,f)$ \textit{contained} in $\Delta(C_1,f)$ -- and the family of objects whose direct sum is taken being induced from a (non-small) family whose index set would be $\Delta(\Fl_{\text{cosquar.}}(\cM),f)$, namely consisting of \textit{all} the square diagrams
\printedpage{11}
% head: A. Grothendieck : Dérivateurs, Ch. XV, Édition des Universités de Montpellier II et Paris VII
% foot: Version 5 novembre 2010
\[
D : \qquad
\begin{tikzcd}
A \arrow[r, "u"] \arrow[d, "i"'] & X \arrow[d, "f"] \\
B \arrow[r, "v"] & Y
\end{tikzcd}
\qquad\qquad \text{with $i$ cosquarable and nothing more ,}
\]
the definition of the $X_D$, $i_D$, $p_D$ associated to $D$ depending only on $D$, and not on the set $C_0$ or $C_1$ of cosquarable arrows which contains $i$. This said, the arrow $u_1(f)$
\msp{page 15}
in (2.5) is none other than the canonical inclusion of a partial amalgamated sum, indexed by $\Delta(C_0,f)$, into the amalgamated sum corresponding to the larger index set. The functoriality of this arrow $u_1(f)$, for variable $f$, ``is immediate'', it follows from the commutativity of the square, associated to a morphism $v=(v_X,v_Y)$ of arrows
\[
\begin{tikzcd}
X \arrow[r, "f"] \arrow[d, "v_X"'] & Y \arrow[d, "v_Y"] \\
X' \arrow[r, "f'"] & Y'\,,
\end{tikzcd}
\tag*{(2.6)}
\]
namely the diagram of maps between index sets
\[
\begin{tikzcd}
\Delta(C_0,f) \arrow[r, hook, "u_1(f)"] \arrow[d, "{\begin{array}{r}\text{functorial dependence}\\ \text{on $f$ of $\Delta(C_0,f)$}\end{array}}"'] & \Delta(C_1,f) \arrow[d, "{\begin{array}{l}\text{functorial dependence}\\ \text{on $f$ of $\Delta(C_1,f)$}\end{array}}"] \\
\Delta(C_0,f') \arrow[r, hook, "u_1(f')"] & \Delta(C_1,f')
\end{tikzcd}
\tag*{(2.7)}
\]
(commutativity which amounts to saying that $\Delta(C_0,f)$ may be regarded as a functor in the pair $(C_0,f)$ as an object of the product category $\underline{\mathfrak{P}_{\text{sm.}}(\Fl_{\text{cosquar.}}(\cM))}\times\uFl(\cM)$ \ldots). The commutative square (2.7) is the ``indices'' part of a commutative square between families of objects $\mathscr{X}(C_0,f)$, $\mathscr{X}(C_0,f')$, $\mathscr{X}(C_1,f)$, $\mathscr{X}(C_1,f')$,
\msp{page 16}
indexed by these sets, namely the $X_D^{C_0}$ and the $X_D^{C_1}$ relative to the diagrams $D$ involving, either the arrow $f$, or the arrow $f'$ \ldots

Now the canonical construction of the ordinal direct system $p_\alpha$, with the $i_\alpha : \mathrm{id}_{\uFl(\cM)}\longrightarrow p_\alpha$, is functorial with respect to the initial datum $p_1$, $i_1$ -- I do not spell out the detail of the functoriality, although I should \ldots

\section*{3\quad Upper bounds for cardinals}

I now suppose that $\cM$ is equipped with a \textit{cardinal filtration}
\[
(\Filt^\pi)_{\pi\geq\pi_0} \qquad\quad (\pi_0,\ \pi \text{ cardinals in } \mathfrak{U})
\tag*{(3.1)}
\]
(SGA 4 I 9.12, p.~149), the $\Filt^\pi$ stable under simple amalgamated sums (whenever they exist), hence also under infinite amalgamated sums indexed by sets
\printedpage{12}
% head: A. Grothendieck : Dérivateurs, Ch. XV, Édition des Universités de Montpellier II et Paris VII
% foot: Version 5 novembre 2010
of indices of cardinality $\leq\pi$. We further suppose that every object of $\Filt^{\pi_0}$ is accessible, and we denote by \circled{$\pi_1$} the smallest cardinal such that $\pi_1\geq\pi_0$ and that the objects of $\Filt^{\pi_0}$ are $\pi_1$-accessible. Denoting, for every cardinal, by $\cM_\pi$ the subset of $\mathrm{Ob}\,\cM$ consisting of the $\pi$-accessible objects, we will then have for every $\pi\geq\pi_1$ (\textit{loc. cit.} 9.16)
\[
\cM_\pi \;\subset\; \Filt^\pi \;\subset\; \cM_{(\pi^{\pi_0})} \qquad\quad (\text{if } \pi\geq\pi_1)\,,
\tag*{(3.2)}
\]

\msp{page 17}

thus if $\pi^{\pi_0}=\pi$, e.g. if $\pi=2^c$ with $c\geq\pi_0$, we will have
\[
\Filt^\pi \;=\; \cM_\pi\,.
\tag*{(3.2 bis)}
\]
(If we admit the generalized continuum hypothesis, we have $\pi^{\pi_0}=\pi$ for every $\pi>\pi_0$, hence $\Filt^\pi=\cM_\pi$ if $\pi\geq\pi_1$ and $\pi>\pi_0$.)

I put, for every cardinal $\pi\geq\pi_0$,
\[
\Fl^\pi(\cM) \;=\; \{\ f\in\Fl(\cM) \mid \text{source and target of $f$ are in } \Filt^\pi\ \}
\tag*{(3.3)}
\]
-- which probably gives a cardinal filtration on $\uFl(\cM)$ (but that is not our concern \ldots). My aim now, it is to bring out conditions on $C_0\subset\Fl_{\text{cosquar.}}(\cM)$ (small subset of $\Fl(\cM)$) and on the cardinal $\pi$, and the ordinal $\alpha$, so that we have
\[
\begin{aligned}
f\in\Fl^\pi(\cM) \;\;\overset{?}{\Longrightarrow}\;\; & X_\alpha(f)\in\Filt^\pi \\
& (\textit{i.e.}\ p_\alpha(f)\in\Fl^\pi(\cM),\ \text{or again } i_\alpha(f)\in\Fl^\pi(\cM))\,.
\end{aligned}
\tag*{(3.4)}
\]
More generally, I propose to give an estimate of the ``cardinality'' of $X_\alpha(f)$, for fixed $C_0$, as a function of $C_0$, of $f$, of $\alpha$, more exactly, as a function of the three cardinals $\pi(C_0)$, $\pi(f)$, $\pi(\alpha)$ defined as
\[
\left\{
\begin{array}{rcl}
\pi(C_0) & = & \text{smallest cardinal such that } C_0\subset\Fl^\pi(\cM) \\[1ex]
\pi(f) & = & \pi(\{f\}) \;=\; \text{smallest cardinal such that } f\in\Fl^\pi(\cM) \\[1ex]
\pi(\alpha) & = & \card\alpha\,.
\end{array}
\right.
\tag*{(3.5)}
\]
Thus I am looking for a cardinal
\msp{page 18}
function of cardinal variables $\varphi$, such that we have
\[
X_\alpha(f) \;\in\; \Filt^\pi \qquad\quad \text{if} \qquad\quad \pi \;=\; \varphi(\underbrace{\pi(C_0)}_{r},\underbrace{\pi(f)}_{s},\underbrace{\pi(\alpha)}_{t})\,.
\tag*{(3.6)}
\]
The crucial estimate will surely concern the case $\alpha=1$. Let us note that $X_1^{C_0}(f)=X_1(f)$ is defined as an amalgamated sum (in general infinite) of objects $X_D$, for a set $\Delta(C_0,f)=\Delta(f)$ of diagrams, namely the commutative diagrams
\[
\begin{tikzcd}
A \arrow[r] \arrow[d, "i"'] & X \arrow[d, "f"] \\
B \arrow[r] & Y
\end{tikzcd}
\]
\printedpage{13}
% head: A. Grothendieck : Dérivateurs, Ch. XV, Édition des Universités de Montpellier II et Paris VII
% foot: Version 5 novembre 2010
with $i\in C_0$. For a fixed $i$, and a fixed $f$, the cardinality of the set of diagrams of this type is bounded above by that of $\mathrm{Hom}(A,X)\times\mathrm{Hom}(B,Y)$. We have $A,B\in\Filt^r$, $X,Y\in\Filt^s$, so we need an upper bound for
\[
\card\mathrm{Hom}_{\cM}(A,X)\,, \qquad\quad A\;\in\;\Filt^r\,, \qquad\quad X\;\in\;\Filt^s\,.
\]

\msp{page 19}

Here are the essential upper bounds in this direction:

\textbf{Proposition 3.1.} \textit{Let $\pi_0$ be the initial cardinal of the cardinal filtration} $(\Filt^\pi)_{\pi\geq\pi_0}$ \textit{under consideration, $\pi_1$ the smallest cardinal $\geq\pi_0$ such that the objects of} $\Filt^{\pi_0}$ \textit{are $\pi_1$-accessible, $C_1=$} card $\underbrace{\widetilde{\uFl^{\pi_1}(\cM)}}_{=\uFl(\cM^{\pi_1})}$\pfn{4}{\textit{cf.} below the notation $\widetilde{\mathcal{C}}$.}. \textit{For every cardinal $\pi\geq\pi_0$, I denote by $\cM^\pi$ the strictly full subcategory of $\cM$ whose set of objects is} $\Filt^\pi$. \textit{Finally, if $\mathcal{C}$ is a category, I denote by $\widetilde{\mathcal{C}}$ the set of isomorphism classes of objects of $\mathcal{C}$.}\pfn{5}{Let us point out that we always have $\card\widetilde{\uFl(\mathcal{C})}\geq\card\widetilde{\mathcal{C}}$ since $\widetilde{\mathcal{C}}\hookrightarrow\widetilde{\uFl(\mathcal{C})}$ via $\mathrm{cl}(X)\longmapsto\mathrm{cl}(\mathrm{id}_X)$.}
% ?? uncertain: straight underlines under "le plus petit" and under "pleine" -- underlined / not underlined (a wavy tilde and an underbrace stand close by)

$1^\circ$) \textit{If $r\geq\sup(\pi_1,C_1)$, we have}
\[
\card\,(\underbrace{\widetilde{\uFl^r(\cM)}}_{=\uFl(\cM^r)}) \;\leq\; 2^{(r^{\pi_0})} \quad\ (=2^r,\ \textit{if } r=2^c,\ \textit{with } c\geq\pi_0)\,.
\tag*{(3.7)}
\]

$2^\circ$) \textit{If $r,s\geq\sup(\pi_1,C_1)$, and if $A\in\Filt^r$, $X\in\Filt^s$, we have}
% (3.8) and (3.9) share one brace on the right
\[
\begin{array}{l}
\text{(3.8)}\\[3ex]
\text{(3.9)}
\end{array}
\qquad\qquad
\begin{array}{r@{\;\;\leq\;\;}l}
\card\mathrm{Hom}_{\cM}(A,X) & s^{(r^{\pi_0})}\\[3ex]
\card\widetilde{\uFl(\cM^r/X)} & s^{(r^{\pi_0})}
\end{array}
\qquad\quad
\left\{
\begin{array}{lll}
= & s^r & \textit{if } r=2^c,\ \textit{with } c\geq\pi_0\,,\\[3ex]
= & s & \textit{if } s=2^{c'},\ \textit{with } c'\geq r^{\pi_0}\,.
\end{array}
\right.
\]

\textbf{N.B.} If we admit the generalized continuum hypothesis, which implies that $\pi^{\pi'}=\pi$ if $\pi>\pi'$ ($\pi$ an infinite cardinal, $\pi'\geq1$), then taking $r$ such that $r>\pi_0$, the preceding formulas give (with moreover $r\geq\sup(\pi_1,C_1)$):
\[
\card\widetilde{\uFl^r(\cM)} \;\;\leq\;\; 2^r\,,
\tag*{(3.7 bis)}
\]
\[
\card\mathrm{Hom}_{\cM}(A,X) \;\leq\; s^r \qquad\quad (\text{if } A\in\Filt^r,\ X\in\Filt^s)\,,
\tag*{(3.8 bis)}
\]
\[
\card\widetilde{\uFl(\cM^r/X)} \;\leq\; s^r \qquad\quad (\text{if } X\in\Filt^s)\,.
\tag*{(3.9 bis)}
\]
Without the continuum hypothesis, these formulas are also valid for $r$ of the form $2^c$, with $c\geq\pi_0$. It is these last formulas which look the most natural. For (3.8 bis) and % [sic] qui on l'air -> qui ont l'air, verb form
(3.9 bis), I point out that the right-hand side $s^r$ is the cardinality of the set $\mathrm{Hom}_{\mathrm{Ens}}(A,X)$, where $A$, $X$ are two sets such that $\card A=r$, $\card X=s$. We therefore cannot hope to improve (3.8 bis)\,! For the

\end{document}
